\section*{Chapitre \ref{ch:b3:behavioural-ecology} --- Écologie comportementale}

\begin{solution}{exo:b3:behavioural-ecology:1}
Mécanisme~: la danseuse transpose l'angle entre la nourriture et le
soleil, mesuré pendant le vol, en un angle à la pesanteur sur le rayon,
et code la distance, tirée du flux optique du trajet, par la durée de la
course. Développement~: largement inné --- des abeilles élevées sans
avoir vu de danse dansent correctement, avec de petits ajustements
appris. Fonction~: elle recrute des congénères vers une source
rentable, ce qui élève les apports de la colonie. Histoire~: dérivée de
mouvements d'intention ritualisés de départ vers la nourriture~; les
abeilles naines dansent encore sur une surface horizontale en pointant
droit vers la source, la forme ancestrale.
\end{solution}

\begin{solution}{exo:b3:behavioural-ecology:2}
Un \omterm{def:b3:behavioural-ecology:tinbergen}{patron moteur fixe} est une séquence stéréotypée déclenchée par un
indice simple, le \omterm{def:b3:behavioural-ecology:tinbergen}{stimulus déclencheur}, et menée à son terme une fois
commencée~: le roulement d'œuf de l'oie, déclenché par un objet en forme
d'œuf~; le picorage du poussin de goéland, déclenché par une tache rouge
sur une forme de bec. Le bâton à trois bandes rouges, picoré plus qu'une
vraie tête, a montré que le mécanisme déclencheur est un filtre accordé
à un trait --- contraste rouge, allongement --- et non une
reconnaissance du parent, et qu'exagérer le trait exagère la réponse.
\end{solution}

\begin{solution}{exo:b3:behavioural-ecology:3}
Quitter une parcelle quand le taux de gain qu'on y obtient est tombé au
taux moyen de l'habitat pris dans son ensemble, trajet compris. Doubler
le temps de trajet abaisse cette moyenne, de sorte que l'animal reste
plus longtemps dans chaque parcelle, épuise chacune plus à fond, et
finit avec un taux global plus bas.
\end{solution}

\begin{solution}{exo:b3:behavioural-ecology:4}
Une SES est une stratégie que, une fois jouée par presque tout le monde,
aucune rare autre ne peut surpasser. La population de Colombes pures
n'en est pas une~: un Faucon isolé gagne chaque affrontement et empoche
$V$ contre le $V/2$ des Colombes. Dans le mélange chaque stratégie gagne
$V(1 - V/C)/2$, moins que le $V/2$ qu'aurait donné la population de
Colombes, parce qu'une fraction des affrontements sont des combats qui
blessent~; personne ne peut faire mieux seul, de sorte que tout le monde
est pris avec la perte.
\end{solution}

\begin{solution}{exo:b3:behavioural-ecology:5}
$V = 6$, $C = 10$~: $p^{*} = 0.6$~; gain $6(1 - 0.6)/2 = 1.2$ pour
chacune, contre $3$ pour la population de Colombes. $V = 12 > C$~: Faucon
est la SES pure, gain $(12 - 10)/2 = 1$, contre $6$ pour les Colombes
--- l'affrontement détruit les cinq sixièmes de la valeur.
\end{solution}

\begin{solution}{exo:b3:behavioural-ecology:6}
$\tau = T_{0}$~: $t^{*} = 1.15\times 2 = \qty{2.3}{min}$, fraction $1 -
\mathrm{e}^{-1.15} = 0.68$, $R = 0.68A/4.3 = 0.16A$ par minute. $\tau =
4T_{0}$~: $t^{*} = 2.0\times 2 = \qty{4}{min}$ (plus précisément $3.9$),
fraction $0.86$, $R = 0.86A/12 = 0.072A$ par minute.
\end{solution}

\begin{solution}{exo:b3:behavioural-ecology:7}
Demi-germains $1/4$~; grand-parent et petit-enfant $1/4$~; cousins
germains $1/8$. Une ouvrière haplodiploïde et son frère~: il est haploïde
et n'a que les gènes de leur mère~; des gènes de l'ouvrière, la moitié
vient de la mère et la moitié de ceux-là sont les copies qu'il a reçues,
donc $r = 1/4$ de son point de vue à elle (du sien à lui c'est $1/2$~:
l'apparentement n'est pas symétrique quand la ploïdie diffère). Reine à
fils~: tous les gènes de celui-ci sont les siens à elle, mais seule la
moitié des siens sont en lui~: $r = 1/2$ de son côté à elle, $1$ du sien.
\end{solution}

\begin{solution}{exo:b3:behavioural-ecology:8}
$n\,r\,b > c$ avec $b = 0.01$, $c = 0.02$~: $n\,r > 2$. Germains, $r
= 1/2$~: $n \ge 5$. Nièces, $r = 1/4$~: $n \ge 9$. Non apparentés~:
jamais --- et c'est pourquoi les mâles qui se dispersent ne crient pas.
\end{solution}

\begin{solution}{exo:b3:behavioural-ecology:9}
Les accouplements doublent par \qty{25}{cm}, la survie baisse d'un
cinquième~: $L = 25$~: $0.5\times 1.2 = 0.6$~; $50$~: $1$~; $75$~:
$2\times 0.8 = 1.6$~; $100$~: $4\times 0.6 = 2.4$~; $125$~:
$8\times 0.4 = 3.2$~; $150$~: $16\times 0.2 =
3.2$. À ces conditions le produit continue de monter bien au-delà des
\qty{50}{cm} naturels, de sorte que le modèle omet des coûts~: la
pénalité de survie d'une longue queue en vol est plus que linéaire, la
réponse des femelles sature, faire pousser les plumes coûte de
l'énergie, et le doublement des nids d'Andersson a été mesuré sur une
saison, non sur une vie. La queue naturelle est là où les coûts réels
mordent.
\end{solution}

\begin{solution}{exo:b3:behavioural-ecology:10}
Riposteur~: contre Faucon $(V - C)/2$, contre Colombe $V/2$, contre
Riposteur $V/2$. Dans une population de Riposteurs purs un Faucon rare
rencontre des Riposteurs, qui ripostent, et empoche $(V - C)/2 < V/2$~:
il ne peut envahir quand $V < C$. Une Colombe rare empoche $V/2$,
exactement le gain des Riposteurs, de sorte que les Colombes sont
neutres et peuvent dériver dedans~; une fois communes, elles laissent
entrer les Faucons (un Faucon empoche $V$ contre une Colombe), et le jeu
à trois stratégies n'a pas de SES simple --- un premier exemple de la
façon dont ajouter une option peut déstabiliser un mélange stable.
\end{solution}

\begin{solution}{exo:b3:behavioural-ecology:11}
Si tout le monde persistait exactement $t_{0}$, un mutant qui
persisterait $t_{0} +
\varepsilon$ gagnerait tous les affrontements pour un surcoût
négligeable, et une population de tels mutants serait battue par
$t_{0} + 2\varepsilon$, et ainsi de suite~: aucun temps fixe n'est
stable, de sorte que la SES doit être un temps aléatoire. Avec une loi
exponentielle la probabilité de persister encore $s$, sachant une
persistance de $t$ jusque-là, vaut $\mathrm{e}^{-s/m}$ quel que soit $t$
--- l'absence de mémoire --- de sorte qu'un adversaire encore en parade
ne dit rien du temps qu'il tiendra encore, et qu'aucune stratégie ne
peut exploiter le temps écoulé.
\end{solution}

\begin{solution}{exo:b3:behavioural-ecology:12}
Les ouvrières valorisent une sœur à $3/4$ et un frère à $1/4$, de sorte
que sous contrôle ouvrier la colonie devrait investir trois fois plus
dans les femelles que dans les mâles, un rapport $3{:}1$ en
investissement~; la reine, également apparentée aux deux, préfère
$1{:}1$. Trivers et Hare (1976) ont trouvé des rapports
d'investissement proches de $3{:}1$ chez des fourmis à reine unique et
monogame~: les ouvrières l'emportent. Les accouplements multiples
abaissent l'apparentement moyen d'une ouvrière à ses sœurs vers $1/4$ à
$1/2$ et poussent l'optimum des ouvrières vers $1{:}1$~; les colonies
d'espèces à reines polyandres investissent bel et bien plus également,
comme l'argument le prédit.
\end{solution}

\begin{solution}{pb:b3:behavioural-ecology:1}
\textbf{1.} $R(t) = A(1 - \mathrm{e}^{-t/T_{0}})/(\tau + t)$~; partir quand
$g'(t) = R(t)$~: $(A/T_{0})\mathrm{e}^{-t/T_{0}} = A(1 - \mathrm{e}^{-t/T_{0}})/
(\tau + t)$.
\textbf{2.} En multipliant par $(\tau + t)\mathrm{e}^{t/T_{0}}/A$~: $(\tau + t)/T_{0}
= \mathrm{e}^{t/T_{0}} - 1$. Avec $\tau = T_{0}$ et $x = t/T_{0}$~: $2 + x =
\mathrm{e}^{x}$~; $x = 1.15$ donne $3.15$ contre $3.16$.
\textbf{3.} $\tau = 4T_{0}$~: $5 + x = \mathrm{e}^{x}$~; $x = 1.94$ donne
$6.94$ contre $6.96$~; $t^{*} \approx 1.94\,T_{0}$, disons $2.0$.
\textbf{4.} Dense~: $t^{*} = \qty{3.5}{min}$, $60(1 - \mathrm{e}^{-1.15}) =
\qty{41}{kJ}$, \qty{68}{\%}, $R = 41/6.5 = \qty{6.4}{kJ/min}$.
Clairsemé~: $t^{*} = \qty{5.8}{min}$, \qty{51}{kJ}, \qty{86}{\%}, $R = 51/17.8 =
\qty{2.9}{kJ/min}$.
\textbf{5.} Cinq minutes~: $g(5) = 60(1 - \mathrm{e}^{-5/3}) = \qty{49}{kJ}$.
Dense~: $49/8 = \qty{6.1}{kJ/min}$ (\qty{96}{\%} de l'optimum)~;
clairsemé~: $49/17 = \qty{2.9}{kJ/min}$ (\qty{99}{\%}). Une règle
grossière perd quelques pour cent --- assez bien pour que la sélection
s'en soit contentée.
\textbf{6.} L'intervalle entre captures s'allonge à mesure que la
parcelle s'épuise, de sorte qu'un temps d'abandon fixe est un taux de
capture marginal fixe --- le critère du théorème, qui est le même dans
chaque parcelle. Elle échoue parce que le bon seuil dépend de
l'habitat~: avec \qty{20}{s} fixes l'animal part trop tôt là où le trajet
est long et trop tard là où il est court, à moins que le temps d'abandon
ne soit lui-même ajusté à l'expérience récente --- ce que font les
animaux en quête.
\textbf{7.} Faucon contre Faucon $(10 - 30)/2 = -10$, Faucon contre
Colombe $10$, Colombe contre Faucon $0$, Colombe contre Colombe $5$~;
$p^{*} = V/C = 1/3$.
\textbf{8.} $W_{H} = 10 - 20p$, $W_{D} = 5(1 - p)$. $p = 0.1$~: $8$ contre
$4.5$, les Faucons se répandent. $p = 1/3$~: $3.33$ chacun, équilibre.
$p = 0.6$~: $-2$ contre $2$, les Colombes se répandent.
\textbf{9.} $3.33$ à la SES contre $5$ pour les Colombes pures~: un tiers
de la valeur de chaque affrontement est perdu à la possibilité du combat.
\textbf{10.} $V = 40 > C$~: Faucon est la SES pure~; chaque affrontement
est un combat. Les années difficiles devraient voir plus de combats et
plus de blessures --- et c'est le cas.
\textbf{11.} Un Faucon qui rencontre un Bourgeois le trouve résident la
moitié du temps (un combat, $(V - C)/2$) et intrus l'autre moitié (une
promenade de santé, $V$)~: en moyenne $(3V - C)/4$. Le Bourgeois empoche
$V/2$ entre les siens. Le Faucon n'envahit que si $(3V - C)/4 > V/2$,
c'est-à-dire $V > C$. Le «~le résident gagne~» du tircis est un
Bourgeois~: une convention à bon compte, stable parce que se battre coûte
plus qu'une tache de soleil ne vaut --- et quand les deux se croient
résidents, tous deux jouent Faucon.
\textbf{12.} La meilleure chose à faire dépend de ce que font tous les
autres, de sorte que la population se fixe là où les stratégies paient
également plutôt qu'au mieux d'un individu.
\textbf{13.} $n\bar r\,b > c$~: $n\bar r > c/b = 3$.
\textbf{14.} Huit germains~: $n\bar r = 4 > 3$, la garde paie. Huit
inconnus~: $0$, jamais. Huit demi-germains~: $2 < 3$, non.
\textbf{15.} $c = 0.01$~: seuil $n\bar r > 1$ --- des demi-germains et
même quelques germains suffisent désormais. L'individu dont le coût est
le plus faible est celui pour qui la règle est satisfaite, de sorte que
les bien nourris montent la garde pendant que les affamés creusent.
\textbf{16.} $k$ sœurs valent $\tfrac{3}{4}k$, $k$ filles
$\tfrac{1}{2}k$~: préférer les sœurs, dans le rapport $3{:}2$.
\textbf{17.} Une sœur prise au hasard est une sœur germaine ($3/4$) avec
la probabilité $1/3$ et une demi-sœur ($1/4$) avec la probabilité
$2/3$~: $\bar r = 0.25 +
0.17 = 0.42 < 0.5$. La préférence s'inverse~: les filles l'emportent
désormais sur les sœurs, de sorte que l'apparentement seul ne peut
expliquer les ouvrières stériles des espèces polyandres --- il faut y
ajouter l'écologie et la police à l'échelle de la colonie.
\textbf{18.} La règle exige seulement que $r$ soit en moyenne plus élevé
parmi les aidés que dans la population~; tout indice corrélé à la
parenté fait l'affaire --- rester près du terrier natal, aider ceux avec
qui l'on a grandi. La règle d'un spermophile peut être «~crier quand on
est près de chez soi~», ce qui sélectionne les apparentés sans en
reconnaître aucun.
\textbf{19.} $W = (L/20)(1 - L^{2}/10^{4})$~: $L = 20$~: $0.96$~; $40$~:
$1.68$~; $60$~: $1.92$~; $80$~: $1.44$.
\textbf{20.} $W' = \tfrac{1}{20}(1 - 3L^{2}/10^{4}) = 0$~: $L^{*} = 100/\sqrt{3}
= \qty{58}{cm}$, $W = 2.89\times\tfrac{2}{3} = 1.92$.
\textbf{21.} $L^{*} = 70/\sqrt{3} = \qty{40}{cm}$~: plus courte.
L'optimum se tient là où le gain marginal en accouplements égale la
perte marginale en survie, de sorte qu'un coût de prédation plus lourd
raccourcit la queue, et que des populations soumises à des prédations
différentes devraient différer d'ornement --- comme le font les guppys
et les xiphos.
\textbf{22.} Les deux. L'emballement prédit que la préférence dépasse le
trait, qui s'arrête à l'optimum de survie~; le handicap prédit que les
femelles préfèrent plus long parce que seul un mâle en forme le porte,
là encore au-delà de ce que la plupart peuvent faire pousser.
L'expérience montre une préférence au-delà de l'optimum et ne tranche
pas entre les explications.
\textbf{23.} Si la queue était gratuite, tout mâle pourrait faire
pousser la plus longue, elle ne dirait rien de sa qualité, et une
préférence pour elle serait exploitée et disparaîtrait. Un coût qui pèse
plus lourd sur un mâle en mauvaise condition rend la longue queue
abordable aux seuls mâles en forme, et la préférence sélectionne alors
la valeur sélective~: c'est le coût qui rend le \omterm{prop:b3:behavioural-ecology:signals}{signal honnête}.
\textbf{24.} Avec un rendement par accouplement huit fois plus grand, la
valeur sélective d'un mâle monte fortement avec les accouplements et
celle d'une femelle presque pas~; les mâles sont sélectionnés pour
rivaliser pour les accouplements --- les affrontements Faucon--Colombe
et leurs armes --- et les femelles, dont chacun des rares descendants
compte, sont sélectionnées pour choisir, ce qui rend les ornements des
mâles rentables. Une seule asymétrie du gain par accouplement produit à
la fois les combats et la parure.
\textbf{25.} $t^{*} \approx \qty{3.5}{min}$ dans l'habitat dense~; $p^{*}
= 1/3$~; le tour de garde paie quand $n\bar r > 3$.
\end{solution}
